\section{CONCLUSION AND FUTURE WORK}

% We propose TAGA, a reactive group-aware navigation model that enables robots to adapt dynamically to human groups using a tangent-based avoidance strategy. By extending an existing simulation framework to incorporate realistic group behaviors, we provide a more comprehensive evaluation of group-aware navigation. Additionally, we introduce Group Collision Rate (GCR) as a new performance metric to quantify group interactions. Our benchmarking results against ORCA, SF, DS-RNN, and AG-RL demonstrate that TAGA significantly reduces group intrusions while maintaining competitive navigation efficiency.

% However, our approach has some limitations that suggest directions for future work. (1) TAGA uses reactive avoidance rather than predictive modeling, limiting its ability to anticipate group movement; incorporating trajectory prediction could improve planning. (2) Evaluations are conducted in simulation, and while the environment captures realistic crowd dynamics, real-world testing is needed for validation. (3) The model assumes stable group structures, but real-world groups may merge or split dynamically, requiring adaptive strategies to handle such transitions. Addressing these challenges will further enhance socially compliant robot navigation in human environments.

% 
\SR{We introduced TAGA, addressing the critical gap between individual-focused and group-aware robot navigation. Our approach uniquely combines tangent-based reactive strategies with existing navigation frameworks, enabling robots to respect human group formations without compromising navigation performance. The substantial reduction in group intrusions across multiple baseline methods demonstrates that the proposed approach significantly enhances social compliance in crowded environments.} %The modular nature of our approach makes it immediately applicable to existing robotic systems without requiring complete redesign or retraining of navigation algorithms.}

However, our approach has some limitations that suggest directions for future work. (1) TAGA uses reactive avoidance rather than predictive modeling, limiting its ability to anticipate group movement; incorporating trajectory prediction could improve planning. (2) Evaluations are conducted in simulation, and while the environment captures realistic crowd dynamics, real-world testing is needed for validation. Addressing these challenges will further enhance socially compliant robot navigation in human environments.
% While TAGA improves group-aware navigation, it relies on reactive avoidance and assumes stable group structures, limiting adaptability to dynamic real-world interactions. Additionally, its evaluation is confined to simulation, requiring validation in real-world settings to ensure robustness in human environments.


% A conclusion section is not required. Although a conclusion may review the main points of the paper, do not replicate the abstract as the conclusion. A conclusion might elaborate on the importance of the work or suggest applications and extensions. 